(1) $E^2 = E$ and $E(1) = P + (1-P) = 1$ are immediate from $P^2 = P$. If $xP = Px$ then the cross-blocks of $x$ vanish ($Px(1-P) = xP(1-P) = 0$ and symmetrically), so $x = PxP + (1-P)x(1-P) = E(x)$; conversely $E(x)$ commutes with $P$ for any $x$; hence $\mathrm{im}\,E = C$ and $E|_C = \mathrm{id}$. The bimodule property follows by moving $c, c'$ through the projectors they commute with. The identification $C = A \oplus A\,ff^*$: $A$ commutes with $P$ (evenness of $f^*f$ in the $H$-factor) and so does $ff^*$, giving $A \oplus A\,ff^* \subseteq C$; and the two diagonal Peirce blocks $PBP \oplus (1-P)B(1-P) = A\,f^*f \oplus A\,ff^*$ span the same $2 \cdot 2^n$-dimensional space, since $a\,f^*f = a - a\,ff^*$. (2) The off-diagonal Peirce blocks contain $Af$ and $Af^*$ ($f = ff^*\cdot f = f \cdot f^*f$ places $f \in (1-P)BP$, symmetrically $f^* \in PB(1-P)$), and the dimensions ($2^n$ each, against Peirce blocks of dimension $2^n$ each in $\Mat_2(A)$) force equality. (3) $E(a) = a$ since $a \in C$. $D(a) = P a P = a\,f^*f$ (as $a$ commutes with $P$ and $P^2 = P$), which is $a$ under the natural corner isomorphism. The Peirce-coordinate reading: the $(2,2)$ entry of $\varphi(a)$ transported to the $E_{11}$ corner is $E_{12}\,a\,E_{21} = f a f^* = \hat{a}\,ff^*$, i.e.\ $\hat{a}$, as computed in Theorem \claimref{transport.rule-1-morita-trivial}. The defect formula is $D(xy) = Px(P + (1-P))yP = D(x)D(y) + Px(1-P)yP$.
